%==============================================================
%  PixelRoot: Hardware-Anchored, Sensor-Level Media Provenance
%  with Real-Time Blockchain Notarization
%
%  Compile:  pdflatex pixelroot  ->  bibtex pixelroot
%            ->  pdflatex pixelroot  ->  pdflatex pixelroot
%  (Requires IEEEtran.cls; available by default on Overleaf.)
%==============================================================
\documentclass[conference]{IEEEtran}
\IEEEoverridecommandlockouts

\usepackage{cite}
\usepackage{amsmath,amssymb,amsfonts}
\usepackage{algorithmic}
\usepackage{graphicx}
\usepackage{textcomp}
\usepackage{xcolor}
\usepackage{booktabs}
\usepackage{pifont}
\usepackage{array}
\usepackage{url}
\usepackage[hidelinks]{hyperref}

\def\BibTeX{{\rm B\kern-.05em{\sc i\kern-.025em b}\kern-.08em
    T\kern-.1667em\lower.7ex\hbox{E}\kern-.125emX}}

\newcommand{\pr}{\textsc{PixelRoot}}

\begin{document}

\title{\pr{}: Hardware-Anchored, Sensor-Level Media Provenance with
Real-Time Blockchain Notarization for Trustworthy Imaging}

\author{\IEEEauthorblockN{Mihir Kavishwar}
\IEEEauthorblockA{\textit{PixelRoot Project} \\
\texttt{hello@mihirkavishwar.com}}}

\maketitle

%==============================================================
\begin{abstract}
The realism of AI-generated ``deepfake'' imagery has outpaced the ability of
society to tell synthetic media from camera-captured reality. Reactive,
detector-based defenses are locked in an adversarial arms race and fail to
generalize to unseen generators, while existing provenance schemes bind trust
to \emph{metadata} that can be stripped, edited, or transplanted. We present
\pr{}, a media-authenticity framework that moves the root of trust from
metadata into the \emph{pixels themselves}. At the instant of capture, a
CMOS image sensor modulates a sparse, pseudorandomly selected set of pixels to
embed a $128$-bit provenance payload---the sensor's factory-burned device
identity, a trusted timestamp, and a quantized capture location---protected by
Reed--Solomon error correction. A \mbox{SHA-256} commitment over the image and
its payload is then notarized to a public blockchain within seconds of the
shutter event, closing the window in which a forgery could be substituted. We
contribute: (i) a structured literature survey unifying five previously
disjoint research threads---deepfake detection, passive sensor forensics
(PRNU), active watermarking, content-provenance standards (C2PA), and
blockchain media integrity---and a gap analysis that motivates a
\emph{dual-binding} design; (ii) a formal architecture and threat model with a
security argument against re-capture, metadata forgery, replay,
compression-laundering, and PRNU fingerprint-copy attacks; (iii) an analysis of
payload survivability under JPEG and \mbox{H.264/HEVC} compression that
clarifies the fragile/semi-fragile trade-off; and (iv) an evaluation plan and a
reference software implementation. We argue that, deployed as an OEM
inter-operability standard analogous to JPEG or \mbox{3GPP}, sensor-level
provenance plus real-time notarization can give legal, journalistic,
third-party (3P), and life-saving applications an affirmative, hardware-rooted
proof of authenticity rather than a probabilistic guess.
\end{abstract}

\begin{IEEEkeywords}
media provenance, deepfakes, blockchain, CMOS image sensor, PRNU, digital
watermarking, content authenticity, C2PA, image forensics, JPEG.
\end{IEEEkeywords}

%==============================================================
\section{Introduction}
\IEEEPARstart{D}{igital} images and videos have become primary evidence in
journalism, courts, insurance, humanitarian monitoring, and everyday social
discourse. That evidentiary role is collapsing. Diffusion models, generative
adversarial networks, and neural rendering now synthesize photorealistic faces,
scenes, and full videos that humans---and automated detectors---routinely fail
to distinguish from reality~\cite{croitoru2024deepfake,verdoliva2020media}. The
harm is no longer hypothetical. Industry incident trackers attribute on the
order of \$1.2--1.3~billion in documented fraud losses to deepfakes in 2025,
roughly tripling year over year, with the overwhelming majority of losses
originating on social-media platforms~\cite{resemble2025threat,surfshark2025fraud};
enterprise surveys report that a majority of organizations experienced
deepfake-related incidents with mean losses exceeding
\$280{,}000~\cite{ironscales2025}. Beyond fraud, non-consensual intimate
imagery and election disinformation erode personal safety and institutional
trust.

The dominant technical response has been \emph{reactive detection}: train a
classifier to spot the statistical fingerprints of a given generator. This
strategy is structurally disadvantaged. Detectors overfit to the artifacts of
the generators in their training set and degrade sharply on unseen
models---recent in-the-wild benchmarks report AUC drops of roughly
$45$--$50\%$ relative to academic
datasets~\cite{chandra2025deepfakeeval,croitoru2024deepfake}. Because every
published detector also serves as a differentiable signal that the next
generation of generators can be trained to defeat, detection is a perpetual
``tug of war'' that the defender is not guaranteed to
win~\cite{tugofwar2024,deepfakevideo2024review}.

An alternative philosophy inverts the problem: instead of trying to prove a
piece of media is \emph{fake}, prove that a piece of media is \emph{real} by
binding it, at creation time, to a trustworthy origin. This is the
\emph{provenance} paradigm now embodied by the C2PA Content Credentials
standard and by secure-capture products such as Truepic and
ProofMode~\cite{c2pa2024spec,truepic2024capture,proofmode2023,google2025c2pa}.
These systems are a major advance, but they share a common structural weakness:
the authenticity claim is carried in a \emph{signed metadata container}
attached to the file. Metadata can be deliberately stripped by platforms,
lost in transcoding, or, worse, the binding can be attacked
(a 2025 camera firmware deployment had to revoke all certificates after a
signing vulnerability~\cite{c2pa2024spec}). If the credential is removed, the
pixels carry no inherent evidence of their origin.

\subsection{The \pr{} thesis}
We argue that a robust root of trust for visual media should live where the
information actually is: in the pixels, anchored to physical hardware, and
witnessed by an immutable public ledger. \pr{} embeds a provenance payload
\emph{into the pixel values} at the moment of capture through controlled
per-pixel sensor gain modulation, and \emph{simultaneously} notarizes a
cryptographic commitment of the image to a blockchain. This yields two
independent bindings:

\begin{itemize}
\item a \textbf{hard binding}---a \mbox{SHA-256} commitment recorded on-chain,
giving an immutable, publicly verifiable timestamp of existence; and
\item a \textbf{soft, in-pixel binding}---a sparse, error-corrected signature
tied to the sensor's factory identity that travels \emph{with the pixels} even
if the file container and its metadata are destroyed.
\end{itemize}

\noindent The two bindings are complementary: the hard binding gives
cryptographic non-repudiation but is destroyed by any re-encoding; the in-pixel
binding tolerates mild processing and survives metadata stripping, but is
statistically softer. Together they cover each other's failure modes.

\subsection{Contributions}
\begin{enumerate}
\item \textbf{A unifying survey and gap analysis} (\S\ref{sec:related})
spanning deepfake generation/detection, passive PRNU forensics, active
watermarking, provenance standards, and blockchain media integrity, showing
that no prior line of work simultaneously achieves (a) a hardware root of
trust, (b) survival of metadata stripping, and (c) public, real-time
notarization.
\item \textbf{A formal architecture} (\S\ref{sec:arch}) for sensor-level
payload embedding, including the $128$-bit metadata layout, PRNG-driven pixel
selection seeded by the payload, Reed--Solomon protection, gain modulation, and
a smart-contract notarization protocol.
\item \textbf{A threat model and security analysis} (\S\ref{sec:threat},
\S\ref{sec:security}) covering re-capture, metadata forgery, replay,
compression-laundering, deepfake substitution, and the PRNU fingerprint-copy
attack, contrasting \pr{} with passive PRNU and metadata-only schemes.
\item \textbf{A compression-robustness analysis} (\S\ref{sec:compression})
relating JPEG quantization and \mbox{H.264/HEVC} coding to payload bit-error
rate, and a discussion of the semantic-authenticity boundary.
\item \textbf{An evaluation plan and reference implementation}
(\S\ref{sec:eval}) grounded in an open prototype, plus a roadmap toward an OEM
standard (\S\ref{sec:discussion}).
\end{enumerate}

%==============================================================
\section{Background}\label{sec:background}

\subsection{The CMOS imaging pipeline}
A modern camera forms an image by integrating photo-generated charge in each
sensor photosite, converting it to a voltage, applying analog gain, and
digitizing it. Manufacturing imperfections make each photosite's response
slightly non-uniform; this multiplicative \emph{photo-response non-uniformity}
(PRNU) is a stable, device-unique noise pattern that has long been used as a
forensic ``sensor fingerprint''~\cite{lukas2006digital,chen2008determining}.
Crucially for \pr{}, many sensors expose programmable per-column or per-region
analog gain and exposure registers, and embed a factory-burned identifier in
one-time-programmable (OTP) memory. \pr{} repurposes these existing controls:
rather than only \emph{reading} the passive PRNU after the fact, it
\emph{actively writes} a chosen, error-corrected signal into the image during
exposure, and ties that signal to the OTP identity.

\subsection{Lossy compression: JPEG and MPEG/H.26x}
Captured media is almost never stored or transmitted losslessly. JPEG applies a
block \mbox{DCT}, quantizes coefficients with a quality-dependent matrix that
preferentially discards high-frequency energy, and entropy-codes the
result~\cite{wallace1992jpeg}. Video codecs (MPEG-2, \mbox{H.264/AVC},
\mbox{HEVC}) add motion-compensated inter-frame prediction, integer transforms,
in-loop deblocking, and group-of-pictures (GOP)
structure~\cite{wiegand2003h264,sullivan2012hevc}. All of these operations are
\emph{lossy} and spatially redistributive, which is exactly why any embedded
signal must be designed to survive quantization---an issue we analyze in
\S\ref{sec:compression}.

\subsection{Cryptographic and blockchain primitives}
\pr{} relies on standard primitives: a collision-resistant hash
(\mbox{SHA-256}~\cite{nist2015sha}) for commitments and PRNG seeding;
Reed--Solomon codes~\cite{reed1960polynomial} for error correction over the
embedded payload; Merkle trees~\cite{merkle1987digital} for batching many
commitments into a single on-chain transaction; and a public, append-only
ledger~\cite{nakamoto2008bitcoin,wood2014ethereum} as an immutable, publicly
auditable timestamping witness. Large media artifacts are referenced---not
stored---on-chain, optionally via content-addressed
storage~\cite{benet2014ipfs}. Hardware attestation of the signing
environment follows established trusted-computing
practice~\cite{trustedcomputing,google2025c2pa}.

%==============================================================
\section{Related Work and Literature Survey}\label{sec:related}
We organize prior art into five threads and then identify the gap \pr{} fills.

\subsection{Deepfake generation and detection}
Surveys by Verdoliva~\cite{verdoliva2020media}, Mirsky and
Lee~\cite{mirsky2021creation}, and Croitoru \emph{et
al.}~\cite{croitoru2024deepfake} chart the rapid progression from face-swap GANs
to diffusion models and neural radiance fields across image, video, audio, and
multimodal content. The recurring empirical finding is poor
\emph{cross-generator generalization}: detectors trained on one family of fakes
collapse on another~\cite{croitoru2024deepfake,deepfakevideo2024review}. The
\mbox{Deepfake-Eval-2024} in-the-wild benchmark quantifies this, reporting
$45$--$50\%$ AUC degradation versus curated academic
sets~\cite{chandra2025deepfakeeval}. The community increasingly frames
generation vs.\ detection as an adversarial co-evolution in which detectors
provide gradient signal to the next generator~\cite{tugofwar2024}. \emph{Limitation:}
detection is reactive, probabilistic, and structurally one step behind; it
cannot offer affirmative, durable proof that a specific artifact is genuine.

\subsection{Passive sensor forensics (PRNU)}
PRNU-based source identification, introduced by Luk\'a\v{s} \emph{et
al.}~\cite{lukas2006digital} and extended to integrity verification by Chen
\emph{et al.}~\cite{chen2008determining}, estimates a camera fingerprint from
residual noise and correlates it (e.g.\ via peak-to-correlation energy) with a
questioned image; recent benchmarks track its accuracy on modern
devices~\cite{prnubench2025}. Broader integrity surveys situate PRNU among
splicing, copy-move, and double-compression
detectors~\cite{korus2017digital,piva2013overview}. \emph{Limitations:} (i)
PRNU is \emph{passive}---it requires a reference fingerprint and many images per
device, and degrades on heavily compressed or low-resolution
media~\cite{prnubench2025}; (ii) it is vulnerable to the
\emph{fingerprint-copy} (PRNU-copy) attack, in which an adversary estimates a
victim's fingerprint from public images and transplants it into a forged image.
The triangle test and its successors detect na\"ive copies but are themselves
defeated by improved, dispersed
copy attacks~\cite{goljan2011defending}---a live cat-and-mouse. \pr{} differs by
\emph{actively} embedding a \emph{chosen, keyed, error-corrected} payload bound
to the OTP identity and a ledger timestamp, rather than relying on an
involuntary, copyable noise pattern.

\subsection{Active authentication: digital watermarking}
Watermarking embeds a signal for authentication or copyright. Schemes are
\emph{robust} (survive processing, for copyright), \emph{fragile} (break on any
edit, for tamper localization), or \emph{semi-fragile} (tolerate benign
operations such as JPEG but break on malicious
edits)~\cite{begum2020digital,fragile2025hybrid}. Classic methods embed in
DCT/DWT/SVD coefficients; recent work uses learned
encoders/decoders~\cite{dlwatermarking2026}. Video-domain schemes embed in
H.264/HEVC transform coefficients or motion vectors and must contend with
quantization and motion-compensation
drift~\cite{videowatermark2018survey,hevcwatermark2020,cstfmark2026}.
\emph{Limitation:} watermarking alone provides no \emph{public, time-anchored}
record and no \emph{hardware} root---a watermark embedded in software can be
embedded by an attacker too. \pr{} treats its in-pixel payload as a
\emph{semi-fragile, hardware-originated} watermark whose trust is bootstrapped
by on-chain notarization and OTP identity.

\subsection{Provenance standards and secure capture}
C2PA Content Credentials define a signed manifest of assertions (capture
device, time, edit history) bound to the asset by a hard hash, verifiable
offline via a certificate chain~\cite{c2pa2024spec}. Hardware-backed signing has
reached consumer devices---Google Pixel signs photos with keys in the Titan~M2
secure element and an on-device timestamp
authority~\cite{google2025c2pa}---and Leica, Sony, and Samsung ship C2PA in
cameras and phones. Truepic's Controlled Capture verifies device integrity via
attestation before signing on-device~\cite{truepic2024capture}; ProofMode
captures sensor metadata, signs with per-capture keys, and registers hashes to
public ledgers (OpenTimestamps/Bitcoin)~\cite{proofmode2023}.
\emph{Limitations:} the binding lives in a \emph{detachable manifest}. If a
platform strips metadata (as many do) or transcodes the asset, the credential
is lost and the bare pixels are again unverifiable; certificate compromise
revokes trust wholesale~\cite{c2pa2024spec}. \pr{} is complementary to
C2PA---it can populate C2PA assertions---but additionally places a redundant
proof \emph{inside the pixels}, surviving manifest loss.

\subsection{Blockchain-based media integrity}
A large body of work anchors media hashes (exact and perceptual) to smart
contracts, storing the bulk content off-chain in
IPFS~\cite{qureshi2021blockchain,benet2014ipfs}. Hasan and Salah propose using
blockchain and smart contracts to trace deepfake video provenance back to a
trusted original~\cite{hasan2019combating}; subsequent systems combine
perceptual hashing, watermarking, and Ethereum/Polygon contracts for video
copyright and
integrity~\cite{videowmblockchain2024,frontiers2025videocopyright}.
\emph{Limitations:} these systems hash media \emph{after} capture, typically in
software at upload time, leaving a tampering window and no hardware root; and
exact hashes are destroyed by any re-encoding while perceptual hashes trade away
cryptographic guarantees. \pr{} narrows the window to the shutter event and adds
a hardware-anchored in-pixel binding the ledger alone cannot provide.

\subsection{Gap analysis}
Table~\ref{tab:compare} summarizes the landscape. No prior approach
simultaneously offers a hardware root of trust, survival of metadata stripping,
\emph{and} public real-time notarization. \pr{} is, to our knowledge, the first
design to combine sensor-level in-pixel embedding with capture-time blockchain
notarization, explicitly targeting that intersection.

\begin{table}[t]
\centering
\caption{Where prior art sits, and the gap \pr{} targets.}
\label{tab:compare}
\renewcommand{\arraystretch}{1.25}
\setlength{\tabcolsep}{3pt}
\begin{tabular}{@{}lccccc@{}}
\toprule
\textbf{Approach} & \textbf{HW} & \textbf{Surv.} & \textbf{Public} & \textbf{Affirm.} & \textbf{Real-} \\
 & \textbf{root} & \textbf{meta-strip} & \textbf{ledger} & \textbf{proof} & \textbf{time} \\
\midrule
Deepfake detectors~\cite{croitoru2024deepfake} & \ding{55} & --- & \ding{55} & \ding{55} & \ding{51} \\
PRNU forensics~\cite{lukas2006digital}        & part. & \ding{51} & \ding{55} & part. & \ding{55} \\
Watermarking~\cite{begum2020digital}          & \ding{55} & \ding{51} & \ding{55} & part. & \ding{51} \\
C2PA / Truepic~\cite{c2pa2024spec,truepic2024capture} & \ding{51} & \ding{55} & opt. & \ding{51} & \ding{51} \\
Blockchain hashing~\cite{hasan2019combating}  & \ding{55} & \ding{55} & \ding{51} & \ding{51} & part. \\
\textbf{\pr{} (this work)}                    & \ding{51} & \ding{51} & \ding{51} & \ding{51} & \ding{51} \\
\bottomrule
\end{tabular}
\\[2pt]
\footnotesize HW root = hardware/sensor root of trust; Surv.\ meta-strip =
survives metadata removal; Affirm.\ proof = affirmative proof of authenticity.
\end{table}

%==============================================================
\section{Threat Model and Design Goals}\label{sec:threat}

\subsection{Assumptions}
We assume: (A1) the sensor's OTP identity and gain registers are written at
manufacture and are not feasibly re-programmable in the field without
specialized equipment; (A2) the capture-time signing/notarization environment is
attested (e.g.\ via a secure element), as in current hardware-backed C2PA
deployments~\cite{google2025c2pa,truepic2024capture}; (A3) the blockchain is
append-only and its history is economically infeasible to rewrite within the
relevant timeframe~\cite{nakamoto2008bitcoin}; (A4) standard cryptographic
primitives are secure~\cite{nist2015sha}.

\subsection{Adversary}
The adversary may: forge or edit file metadata; generate a deepfake or
Photoshop a real image; re-encode/transcode/compress media; re-photograph a
screen (analog ``re-capture''); attempt to estimate and transplant a victim's
sensor fingerprint (PRNU-copy~\cite{goljan2011defending}); and replay or
back-date claims. The adversary cannot economically clone a specific physical
CMOS die or rewrite ledger history.

\subsection{Design goals}
\textbf{G1 Affirmative authenticity:} produce verifiable evidence that media
originated from a genuine sensor at a specific time/place.
\textbf{G2 Hardware root:} bind the proof to physical sensor identity, not to
software-only secrets.
\textbf{G3 Metadata-independence:} retain evidentiary value even if the file
container/metadata is stripped.
\textbf{G4 Public verifiability:} allow any third party (platform, court,
viewer) to verify without contacting the originator or possessing the camera.
\textbf{G5 Minimal capture latency:} notarize within seconds of the shutter.
\textbf{G6 Imperceptibility:} embedding must not visibly degrade the image.
\textbf{G7 Standardizability:} be deployable as a cross-OEM interoperability
standard.

%==============================================================
\section{The \pr{} Architecture}\label{sec:arch}

\pr{} comprises a capture-time \emph{embed-and-notarize} pipeline and a
\emph{verification} pipeline (Fig.~\ref{fig:arch} conceptually; described
below).

\subsection{Provenance payload}
The payload $\mathbf{m}$ is a fixed-width record (\pr{}~v2: $128$ bits):
\begin{equation}
\mathbf{m} = \underbrace{\mathrm{ID}}_{48} \,\|\,
             \underbrace{t}_{32} \,\|\,
             \underbrace{(\phi,\lambda)}_{48},
\end{equation}
where $\mathrm{ID}$ is the sensor's OTP device identifier (modeled here as a
MAC-style $48$-bit value), $t$ is a Unix timestamp from an attested
RTC/NTP source, and $(\phi,\lambda)$ is a quantized latitude/longitude
($24$ bits per coordinate, $\approx\!10^{-4}$ degree resolution). Location may be
omitted or salted for privacy (\S\ref{sec:discussion}).

\subsection{Keyed pixel selection}
To make the embedding location unpredictable and content-spreading, the
$n_b$ payload-carrying pixels are chosen by a PRNG seeded by the payload itself:
\begin{equation}
s = \mathrm{trunc}_{32}\big(\mathrm{SHA\text{-}256}(\mathbf{m})\big),
\end{equation}
and positions $\{(r_i,c_i)\}_{i=1}^{n_b}$ are drawn without repetition from the
$H\times W$ sensor grid using a PRNG keyed by $s$. Because the seed is derived
from $\mathbf{m}$, a verifier who is told the claimed $\mathbf{m}$ can regenerate
the exact positions; an attacker who does not know $\mathbf{m}$ cannot localize
the carriers. (A production deployment additionally mixes in an OEM secret key
$k$, i.e.\ $s=\mathrm{trunc}_{32}(\mathrm{SHA\text{-}256}(k\,\|\,\mathbf{m}))$,
so that carriers cannot be enumerated even given $\mathbf{m}$.)

\subsection{Error-corrected encoding}
The $128$ payload bits are expanded with a Reed--Solomon
code~\cite{reed1960polynomial} into $n_b=128+r$ coded bits ($r\approx 30$
parity bits in the prototype), tolerating burst errors from compression and
mild edits. Let $b_i\in\{0,1\}$ be the $i$-th coded bit.

\subsection{Sensor-level gain modulation (embedding)}
During exposure, each selected photosite's analog gain is nudged
multiplicatively:
\begin{equation}
g_i = 1 + \alpha\, b_i, \qquad \alpha \approx 0.02,
\end{equation}
so a ``1'' raises the pixel response by $\approx\!2\%$ ($\approx\!0.01$~EV) and a
``0'' is left nominal. The perturbation is below the perceptual threshold and is
spread across $\sim\!10^{2}$ pixels out of millions (goal~G6). This is an
\emph{active}, hardware-rooted, semi-fragile watermark
(cf.\ \S\ref{sec:related}-C), distinct from the involuntary PRNU pattern that
passive forensics merely observes.

\subsection{Commitment and notarization}
Immediately after readout, the device computes a hard binding
\begin{equation}
h = \mathrm{SHA\text{-}256}\big(\,I \,\|\, \mathbf{m}\,\big),
\end{equation}
over the (lossless) image $I$ and payload, then submits $h$ (optionally with a
device pseudonym and coarse time/location) to a notarization smart contract.
For scale, many commitments are aggregated into a Merkle
tree~\cite{merkle1987digital} and only the root is written on-chain, amortizing
gas cost; the contract emits an event with a block timestamp that serves as the
public ``proof of existence''~\cite{proofmode2023}. If connectivity is briefly
unavailable, $h$ is queued in secure storage and submitted on reconnect, with
the latency recorded.

\begin{algorithm}[t]
\caption{\pr{} capture-time embed-and-notarize}
\label{alg:capture}
\begin{algorithmic}[1]
\STATE $\mathrm{ID}\!\gets\!\mathrm{readOTP}();\;
        t\!\gets\!\mathrm{attestedClock}();\;
        (\phi,\lambda)\!\gets\!\mathrm{gps}()$
\STATE $\mathbf{m}\gets \mathrm{ID}\,\|\,t\,\|\,(\phi,\lambda)$
\STATE $s\gets \mathrm{trunc}_{32}(\mathrm{SHA\text{-}256}(k\,\|\,\mathbf{m}))$
\STATE $\mathbf{b}\gets \mathrm{ReedSolomonEncode}(\mathbf{m})$
\STATE $P\gets \mathrm{PRNGselect}(s,\,|\mathbf{b}|,\,H,\,W)$
       \quad\COMMENT{carrier pixels}
\FORALL{$(r_i,c_i)\in P$}
   \STATE $\mathrm{setGain}(r_i,c_i,\,1+\alpha\, b_i)$
\ENDFOR
\STATE $I\gets \mathrm{expose\&readout}()$
\STATE $h\gets \mathrm{SHA\text{-}256}(I\,\|\,\mathbf{m})$
\STATE $\mathrm{notarize}(h)$ \quad\COMMENT{Merkle-batched on-chain commit}
\STATE \textbf{store} $I$ (lossless preferred) with $\mathbf{m}$ in C2PA manifest
\end{algorithmic}
\end{algorithm}

\subsection{Verification}
A verifier with a questioned image $I'$ and a claimed payload $\mathbf{m}'$
(from the manifest, the on-chain record, or asserted by a party) performs two
independent checks:

\noindent\textbf{(V1) Hard binding / ledger check.} Recompute
$h'=\mathrm{SHA\text{-}256}(I'\|\mathbf{m}')$ and confirm $h'$ appears in the
ledger (directly or via a Merkle proof to the recorded root). A match proves
$I'$ existed, bit-for-bit, at the recorded block time. This is conclusive but
brittle: any re-encoding of $I'$ breaks it.

\noindent\textbf{(V2) In-pixel binding check.} Regenerate $s$ and the carrier
set $P$ from $\mathbf{m}'$ (and $k$), read back the relative gain at each
carrier to recover noisy coded bits $\hat{\mathbf{b}}$, Reed--Solomon-decode to
$\hat{\mathbf{m}}$, and test $\hat{\mathbf{m}}\stackrel{?}{=}\mathbf{m}'$. The
\emph{bit-error rate} (BER) and the decode success constitute a soft authenticity
score that degrades gracefully under compression
(\S\ref{sec:compression}). Because $P$ is keyed and spread, an attacker who edits
or splices the image without knowing $P$ corrupts the signature in a detectable
way.

\noindent A piece of media is accepted as authentic at the strongest level
supported by the evidence: V1$\wedge$V2 (pristine original), V2 only
(transcoded but pixel-consistent), or V1 only (exact copy, manifest intact).

%==============================================================
\section{Security Analysis}\label{sec:security}

\textbf{Deepfake / synthetic substitution.} A fully AI-generated image was never
exposed on a genuine sensor; it has no valid on-chain commitment (V1 fails) and
cannot reproduce a payload-consistent, keyed in-pixel signature tied to a real
OTP identity (V2 fails). The attacker would need to either forge a ledger entry
(infeasible, A3) or embed a valid signature without the OEM key and sensor
(infeasible, A1--A2).

\textbf{Metadata forgery / stripping.} Editing or removing file metadata cannot
manufacture a matching ledger commitment, and cannot create the in-pixel binding
(goal~G3). Stripping metadata downgrades a credential-based scheme to
``unverifiable,'' but \pr{} still recovers $\mathbf{m}$ from the pixels via V2.

\textbf{Re-capture (photo-of-a-screen).} Re-photographing authentic media on a
genuine \pr{} sensor produces a \emph{new} payload (new $\mathrm{ID}$, $t$,
location) and a \emph{new} commitment; it does not reproduce the original's
ledger entry. The semantic content is duplicated but the provenance correctly
reflects the re-capture event and time, exposing back-dating attempts. Liveness
heuristics (moir\'e, depth) as used by Truepic~\cite{truepic2024capture} can
augment this.

\textbf{Replay / back-dating.} Because notarization is bound to a block
timestamp, an adversary cannot claim an artifact predates its true creation: the
earliest provable existence is the block time. This is the property reactive
detection fundamentally lacks.

\textbf{PRNU fingerprint-copy.} The strongest hardware-targeted attack against
passive forensics estimates a victim's PRNU from public images and transplants
it~\cite{goljan2011defending}. \pr{} resists this for three reasons: (i) the
\pr{} signature is an \emph{active, keyed} code, not the involuntary PRNU, so
copying the passive noise does not reproduce the carrier bits without the OEM key
$k$; (ii) even a perfectly transplanted signature must still match an on-chain
commitment made at capture time by a genuine device (V1/A3); and (iii) the
triangle-test family of countermeasures remains available as a forensic backstop
on the passive PRNU channel. Thus \pr{} converts a soft, copyable fingerprint
into a key-protected, ledger-witnessed credential.

\textbf{Compression-laundering.} An attacker may aggressively re-encode media to
erase the signal. Heavy compression indeed breaks V1 and, past a threshold, V2
(\S\ref{sec:compression}); but this yields a \emph{negative} result
(``cannot verify''), not a \emph{false positive} (``verified authentic''). \pr{}
fails safe: it never certifies a manipulated artifact as genuine; it only
declines to certify.

%==============================================================
\section{Robustness Under Compression}\label{sec:compression}

The central engineering tension is that the hard binding (V1) is \emph{fragile}
by design---ideal for a pristine original---while real-world distribution
demands \emph{semi-fragile} survivability (V2). We therefore analyze the in-pixel
channel under lossy coding.

\subsection{JPEG}
JPEG's quantization matrix scales with the quality factor $Q$ and most
aggressively discards high-frequency DCT energy~\cite{wallace1992jpeg}. A
$\sim\!2\%$ gain perturbation at isolated pixels manifests partly as
high-frequency content and is therefore attenuated as $Q$ drops. To survive, the
in-pixel channel should: (i) place carriers to influence \emph{mid-frequency}
DCT coefficients of their $8\!\times\!8$ blocks rather than single pixels, as in
semi-fragile DCT watermarking~\cite{begum2020digital,fragile2025hybrid}; (ii)
rely on Reed--Solomon parity to absorb the residual bit errors; and (iii) use a
soft, correlation-based read-back rather than a hard per-pixel threshold. We
expect a monotonic BER--$Q$ relationship: near-zero BER at $Q\!\gtrsim\!90$,
graceful rise through mid quality, and decode failure (safe ``cannot verify'')
at low $Q$. The exact crossover is the key quantity to characterize empirically
(\S\ref{sec:eval}).

\subsection{Video: MPEG / H.264 / HEVC}
Video adds motion-compensated prediction, integer transforms, in-loop
deblocking, and GOP structure~\cite{wiegand2003h264,sullivan2012hevc}. The video
watermarking literature shows that embedding survives best in mid-frequency
transform coefficients and that I-frame perturbations propagate (``drift'')
through a GOP, so embedding is often confined to portions of the
GOP~\cite{videowatermark2018survey,hevcwatermark2020,cstfmark2026}. For \pr{},
the practical design is: embed the per-frame (or per-keyframe) payload at
capture, and commit a Merkle root over keyframe hashes so that
\emph{frame-level} verification localizes tampering. Learned, codec-aware
embedding has recently been shown to survive non-differentiable
\mbox{H.264}~\cite{cstfmark2026} and is a natural upgrade path.

\subsection{The semantic-authenticity boundary}
A subtle limitation deserves emphasis. A benign platform filter may alter
carrier pixels enough to fail V2 even though the image's \emph{meaning} is
unchanged---a false alarm---while a malicious edit might, in principle, preserve
carriers. \pr{} therefore certifies \emph{pixel-level integrity relative to a
notarized original}, not semantic equivalence. We argue this is the right
primitive: it cleanly separates ``bit-for-bit original'' (V1), ``faithfully
transcoded'' (V2), and ``cannot establish provenance,'' and it composes with
perceptual-hash and content-similarity layers~\cite{videowmblockchain2024} for
softer ``same-scene'' judgments.

%==============================================================
\section{Applications: Human-in-the-Loop Trust}\label{sec:apps}

\pr{} is designed so that a human or institutional verifier remains in the loop
with hardware-rooted evidence rather than a model's opinion.

\textbf{Legal and evidentiary.} Chain-of-custody for photographic/video evidence
benefits from an immutable, independently timestamped commitment made at
capture, reducing disputes over authenticity and time of
creation~\cite{korus2017digital}.

\textbf{Journalism and third-party (3P) verification.} Platforms and fact-checkers
can verify an asset against the public ledger without contacting the
source---valuable for conflict and human-rights documentation, where
ProofMode-style field capture has already been deployed~\cite{proofmode2023}.

\textbf{Life-saving / high-assurance.} Insurance inspections, telemedicine
imagery, infrastructure monitoring, and emergency reporting need affirmative
proof that an image is a real, current capture of a real scene---exactly the
``prevent rather than detect'' posture argued by secure-capture
vendors~\cite{truepic2024capture}.

\textbf{Platform-scale screening.} Social platforms can perform an
$O(1)$ ledger lookup at upload to label assets ``camera-original,''
``transcoded-but-consistent,'' or ``unverified,'' shifting the default from
``trust until debunked'' to ``label by provenance.''

%==============================================================
\section{Evaluation Plan and Reference Implementation}\label{sec:eval}

\subsection{Reference implementation}
We provide an open software prototype of the \pr{} core (payload layout, SHA-256
seeding, PRNG carrier selection, gain-value generation, image hashing, and a
register/verify service with simulated notarization). It implements both a
legacy \mbox{v1} ($120$-bit) and the current \mbox{v2} ($128$-bit) payload, the
latter fixing a coordinate-packing precision bug. The prototype lets us validate
the encode/verify round-trip and the metadata$\rightarrow$seed$\rightarrow$
position determinism end-to-end before silicon.

\subsection{Proposed experiments}
\begin{enumerate}
\item \textbf{Imperceptibility (G6):} PSNR/SSIM of embedded vs.\ nominal capture
across $\alpha\in[0.01,0.05]$ and carrier counts.
\item \textbf{Compression robustness:} BER and Reed--Solomon decode-success of
V2 vs.\ JPEG quality $Q\in[50,100]$ and vs.\ H.264/HEVC at varied QP/bitrate;
report the safe ``cannot verify'' crossover.
\item \textbf{Latency (G5):} shutter$\rightarrow$hash and
hash$\rightarrow$on-chain confirmation on testnets, including offline-queue
behavior.
\item \textbf{Security:} empirical false-accept under deepfake substitution,
metadata forgery, re-capture, and a simulated PRNU-copy attack; confirm zero
false-accepts (fail-safe property).
\item \textbf{Capacity/ECC trade-off:} payload size vs.\ parity vs.\ robustness.
\item \textbf{Baselines:} compare against metadata-only C2PA (manifest stripped)
and software upload-time hashing on survival of metadata removal and transcode.
\end{enumerate}

%==============================================================
\section{Discussion: Limitations, Privacy, Standardization}\label{sec:discussion}

\textbf{Hardware dependency.} \pr{}'s strongest guarantees require sensor and
secure-element support; legacy devices can run a software-only, attested variant
that still notarizes a capture-time hash (\`a la C2PA/ProofMode) but without the
hardware-rooted in-pixel binding.

\textbf{Privacy.} $\mathrm{ID}$, time, and location can constitute PII. We
recommend on-chain pseudonyms, optional/omittable and salted location, and
selective disclosure so that the public record proves \emph{existence and
integrity} without revealing the photographer's identity or precise whereabouts.

\textbf{Failure-safe semantics.} As emphasized, \pr{} declines rather than
falsely certifies under heavy laundering; downstream policy must treat
``unverified'' as ``unknown,'' not ``fake.''

\textbf{Toward an OEM standard.} The historical lesson of JPEG~\cite{wallace1992jpeg}
and 3GPP is that interoperability emerges from a shared standard, not point
products. We propose \pr{} as a cross-OEM capture-provenance profile that (a)
fixes payload layout, ECC, and the keyed selection function; (b) reuses C2PA
manifests as the transport for assertions; and (c) standardizes the notarization
contract interface, so any platform can verify any vendor's media. Consumer
hardware-backed C2PA signing~\cite{google2025c2pa,truepic2024capture} shows the
ecosystem is ready for the signing half; \pr{} adds the in-pixel, ledger-witnessed
half.

%==============================================================
\section{Conclusion and Future Work}\label{sec:conclusion}
Reactive deepfake detection cannot win a generator-vs-detector arms race, and
metadata-only provenance breaks the moment a credential is stripped. \pr{}
relocates the root of trust into the pixels and onto a public ledger: a
hardware-anchored, error-corrected, keyed signature embedded by the CMOS sensor
at capture, plus a real-time blockchain commitment. This dual binding gives
affirmative, publicly verifiable, hardware-rooted proof of authenticity that
survives metadata loss and fails safe under laundering. Future work includes
silicon-level validation on sensors with per-region gain control, codec-aware
learned embedding for video~\cite{cstfmark2026}, privacy-preserving selective
disclosure, formal modeling of the V2 soft-authenticity score, and pursuing
\pr{} as an open cross-OEM standard. We view this as a foundational---if not
all-encompassing---step toward driving deepfake-driven disinformation toward
practical irrelevance for media that matters.

%==============================================================
\section*{Reproducibility and Disclosure}
The reference implementation is part of the open \pr{} project. The on-chain
notarization in the current prototype is \emph{simulated} and labeled as such;
results reported as ``planned'' are part of the evaluation roadmap
(\S\ref{sec:eval}).

\bibliographystyle{IEEEtran}
\bibliography{references}

\end{document}
